%----------------------------------------------------------------------------
\chapter{Discussion}
\label{chap:Discussion}
%----------------------------------------------------------------------------

%----------------------------------------------------------------------------
\section{Synthesis across axes}
\label{sec:Synthesis}
%----------------------------------------------------------------------------
\begin{itemize}
    \item Efficiency: which solver dominates the equal-wall-clock frontier, and in which scenario?
    \item Robustness: where do the stability envelopes diverge, and does the divergence match the primal/dual prediction?
    \item Realism: which solver's converged state is closer to Newton at matched compute?
    \item Flexibility: how does the ease-of-extension cost compare against the per-frame cost benefit?
\end{itemize}

%----------------------------------------------------------------------------
\section{Verdicts across the four axes}
\label{sec:Verdicts}
%----------------------------------------------------------------------------

\paragraph{Efficiency.}
On the equal-wall-clock frontier (\autoref{ssec:ErrorVsMs}) neither solver
dominates outright; the winner depends on how hard the per-step problem is. A
single VBD iteration always lands closer to the Newton reference than a single
XPBD iteration, but it costs more. When the time step is large --- few substeps,
a stiff sub-problem --- VBD reaches a given accuracy for less wall-clock, because
it spends its budget on iterations that carry second-order information about the
energy. As the step shrinks, each sub-problem becomes easy enough for both
solvers to converge in one or two iterations, and XPBD's cheaper per-iteration
cost lets it take the frontier. The scaling test (\autoref{ssec:NodeCountVsFPS})
tells the same story at fixed settings: with early stopping,
$VBD_{sub4\times it_{\max}4}$ beats XPBD on the easy, low-resolution scenes and
converges toward it as the mesh grows, while the cheaper $VBD_{sub4\times it2}$
configuration outperforms $XPBD_{sub4\times it4}$ across the board. These figures
are single-threaded and therefore understate VBD's standing
(\autoref{ssec:ThreatSingleThread}).

\paragraph{Robustness.}
This is the axis where the primal/dual distinction (\autoref{sec:PrimalDual})
makes a concrete prediction, and the two stress tests confirm it directly. Under
a $1000\times$ mass ratio (\autoref{ssec:ChainWithHeavyWeight}) XPBD --- the dual
solver --- fails to reach the Newton reference within 100 iterations, while VBD
approaches it: a dual update propagates a mass disparity only one shared particle
per Gauss--Seidel sweep. Under a $1{:}10000$ stiffness ratio
(\autoref{ssec:StiffnessRatioChain}) the roles reverse: VBD --- the primal
solver --- stalls while XPBD matches the reference almost exactly, because the
per-vertex Hessian block collects every incident spring stiffness and becomes
ill-conditioned. The self-collision drop (\autoref{ssec:SelfCollisionDrop}) is
the same VBD weakness in a practical setting: my 12-spring cloth model carries
exactly the wide stiffness ratios that hurt VBD, so the cloth was unstable unless
I raised vertex mass by two-to-three orders of magnitude relative to the largest
stiffness. XPBD tolerated these constraints, including near-rigid ones, without
hand-tuning. Overall XPBD is the more robust solver for stiff cloth, with the
caveat that this failure is specific to canonical 2024 VBD and is reportedly
addressed by AVBD (\autoref{ssec:ThreatAVBD}).

\paragraph{Physical realism.}
At matched, sufficient iteration counts the two solvers are essentially
indistinguishable from Newton and from each other. The harmonic oscillator
(\autoref{ssec:HarmonicOscillator}) converges in a single iteration for both,
leaving only the damping intrinsic to implicit Euler. The energy traces
(\autoref{ssec:EnergyTracking}) follow one shared exponential decay; at four
iterations both sit on the reference, and the only ordering visible at very low
budgets is that one VBD iteration falls between one and two XPBD iterations ---
the same per-iteration-quality story as the efficiency axis, since the residual
oscillations are convergence artefacts rather than a physical difference.
Realism is therefore a question of how much iteration budget each solver needs to
reach a converged state, not of the converged state itself.

\paragraph{Flexibility.}
Here the comparison favours each solver on a different sub-question. Both make
the converged shape independent of iteration count
(\autoref{ssec:IterIndepStiff}), so neither couples material stiffness to solver
budget. For adding new constraints (\autoref{ssec:NewFunctionality}), XPBD was
easier to extend with the scalar dihedral-bending constraint: it needs only the
first derivative and a single extra solve call, whereas VBD needs the second
derivative and stores each hinge once per incident vertex, recomputing its
Hessian block four times an iteration. This advantage is constraint-shaped rather
than absolute --- for tensor-valued energies such as StVK it would be XPBD doing
the extra work. On domain coverage (\autoref{ssec:DomainCoverage}) XPBD's decade
of deployment across cloth, fluids, rigid bodies and hair is decisive against a
2024 method still largely confined to elastic bodies.

\paragraph{Joint reading.}
No solver wins on every axis, so the honest recommendation is scenario-based.
VBD is the better choice when the per-step problem is hard --- large time steps,
heavy mass ratios --- and when the target is parallel hardware that suits its
per-vertex blocks. XPBD is the better choice when constraints are stiff or
near-rigid, when robustness out of the box matters more than peak accuracy per
millisecond, or when the breadth of an established ecosystem is an asset. The
primal/dual lens earns its place here not as a verdict on which solver is better,
but as the thing that predicted --- ahead of the measurements --- exactly which
scenario would break which solver.

\paragraph{Additional remarks.}
Based on my experiments, there is a gread difference in the failure of the two solvers.
When XPBD fails, it becomes really stretchy but still looks plausible. VBD however either shows noticable artifacts at failure of explodes.

%----------------------------------------------------------------------------
\section{Limitations and Threats to validity}
\label{sec:ThreatsValidity}
%----------------------------------------------------------------------------

\subsection{Single-threaded execution}
\label{ssec:ThreatSingleThread}
\begin{itemize}
    \item All measurements run single-threaded on CPU.
    \item Because VBD parallelises better than XPBD (\autoref{sec:Parallelisation}), the ms/frame figures are biased against VBD: they understate the advantage it would show in a production, multi-threaded deployment.
    \item This is a threat to the external validity of the quantitative efficiency claims only; it leaves the qualitative claims about convergence quality and stability envelopes untouched, since those do not depend on how the work is scheduled across cores.
    \item Lifting this limitation --- the multithreaded and GPU comparison --- is the principal direction of \autoref{ssec:FWParallel}.
\end{itemize}

\subsection{Mass-spring discretisation only}
\label{ssec:ThreatMassSpring}
\begin{itemize}
    \item I use distance-spring energies throughout, including for bending. Dihedral bending constraints were implemented (with LLM assistance) but I deliberately did not run the evaluation on this newer model.
    \item This choice is not neutral between the solvers: VBD's original paper uses an energy-based model (StVK), and such a discretisation would likely favour VBD to a similar degree that my mass-spring model favours XPBD.
    \item The results are therefore not necessarily transferable to other cloth models. The primal/dual framing still carries over --- both XPBD and VBD admit arbitrary elastic energies --- but the specific stiffness ratios encountered, and hence the magnitude of the stiffness-ratio and per-iteration-convergence effects, may differ.
    \item Re-running the comparison on energy-based and alternative-connectivity models is taken up as future work in \autoref{ssec:FWFEM}.
\end{itemize}

\subsection{Gauss--Seidel order dependence}
\label{ssec:ThreatOrder}
\begin{itemize}
    \item Both solvers are order-sensitive (\autoref{sec:GSJacobi}).
    \item I fix a single deterministic order; an adversarial reordering could in principle shift results.
    \item I do not investigate this, and it is a real threat to reproducibility on other implementations.
\end{itemize}

\subsection{Augmented Vertex Block Descent}
\label{ssec:ThreatAVBD}
\begin{itemize}
    \item I do not implement AVBD or any other post-2024 VBD extension; the comparison is restricted to canonical 2024 VBD.
    \item AVBD fixes several weaknesses of canonical VBD (notably wide stiffness ratios and near-infinite stiffness), so the robustness and stiffness-ratio findings that go against VBD should not be read as applying to AVBD.
    \item Whether those findings are intrinsic to the primal approach or specific to the 2024 implementation is exactly the question \autoref{ssec:FWAVBD} proposes to settle.
\end{itemize}